%  Draft 2026-11-16: full Discussion rewrite (author request: "review the entire discussion to write like a
%  high-profile journal"; Discussion must interpret results, not restate them). Not yet in sections/discussion.tex.
%  Order: summary; joint versus post hoc evidence (MethPhaser); roles of the evidence classes; withholding;
%  benchmarks; limitations; outlook and closing. Every number is taken from the current Results.

Long reads observe several classes of variation on the same molecules, but phasers have used these classes separately or one after another. \toolname places SNVs, indels, SVs and 5mC in one weighted graph and withholds phase where the local graph predicts a switch. On two sequencing platforms, at 10--60$\times$ coverage and in a second individual, this design gave lower SNV switch error rates than read-based and methylation-aware phasers, lower SV phasing error rates than \whatshap, and longer blocks when all four classes were phased together. The cost was a reduction in SNV recall of at most 1 percentage point. These results suggest that, in long-read phasing, how a method combines its evidence and which assignments it reports affect accuracy more than the amount of evidence it adds.

Where methylation enters the analysis determines which errors it can prevent. Earlier methods used methylation or SVs to connect blocks in regions with few heterozygous SNVs\cite{akbari2021nanomethphase,fu2024methphaser,akbari2023parentoforigin,lin2022longphase,holt2024hiphase}; \toolname instead weighs them together with SNVs and indels. The comparison with MethPhaser shows why this difference matters. MethPhaser joins the blocks of a fixed SNV phasing and cannot revise phase within them, so one incorrect join places a whole block on the opposite haplotype. Its blocks were longer but carried a larger fraction of SNVs on the wrong haplotype, an error that the switch error rate counts only once (Fig.~\ref{fig:methylation}d,e and Supplementary Fig.~1e). Joint weighing may reject joins that the combined evidence does not support, at the cost of shorter blocks. Block N50 alone therefore overstates the benefit of methylation, and evaluations of methods that join blocks should report Hamming distance together with switch error rates.

The evidence classes had different roles, and the number of sites that a class added did not predict its effect. 5mC added about twice as many phased sites as indels but joined 4.4--9.5-fold fewer blocks per added site, with little change in switch error rates. Because a 5mC site is kept only when its methylation co-segregates with a neighbouring heterozygous variant, many such sites may lie inside blocks that sequence variants already connect. The blocks joined by indels and by 5mC were nearly additive, which indicates that the two classes bridge different gaps. Indels joined the most blocks but added switch errors with each join. \whatshap showed a similar trade-off, which suggests that this cost lies in indel calls and their representation, particularly in tandem repeats, rather than in one phasing model; more accurate indel calls\cite{zheng2026accelerated,kolesnikov2024local} could lengthen blocks without it. SVs barely changed block length, as expected from a design in which an SV never connects two blocks on its own. Their benefit lies in their own phase: \toolname phased SVs with lower error rates than \whatshap on both platforms and with the lowest rates for SVs of at least 1~kb, whereas HiPhase was more accurate for shorter SVs on HiFi data.

Withholding phase changes what a phaser reports, not only how accurately it phases. The calls that \toolname left unphased rarely matched the benchmark, and it phased 29--48\% fewer genotype errors than \whatshap; each such error, once phased, places an incorrect allele on a haplotype. Benchmark SNVs that \toolname left unphased were phased incorrectly by \whatshap three to six times as often as other SNVs. Withholding therefore concentrated on the calls and sites where phase was least reliable. Because withholding phase and splitting blocks change which variant pairs are scored, switch error rates are interpretable only together with recall, precision and contiguity, and phasers should report which calls they phase. Most unphased calls were absent from the benchmark and often lay in high-depth or structurally complex regions. Their status is therefore unresolved, not shown to be wrong.

Benchmark choice changed the apparent separation between phasers. Under GIAB v4.2.1, most switch errors of all three phasers fell at the same positions and the differences narrowed with coverage; under v5.0q, few errors were shared and the phasers separated more clearly, particularly in segmental duplications. Errors shared by phasers with different algorithms, including \whatshap and \hapcut, which have no parameters trained on either benchmark, are consistent with errors in the phase of the earlier benchmark. Shared reads, alignments and variant calls could also produce correlated errors, and the reads at these sites have not been examined. Much of the genome-wide difference between phasers also arose outside the regions of both benchmarks, where the reference phase is least certain, and v5.0q was released to evaluate variant calls rather than phase\cite{olson2026giabv5}. Phasing benchmarks should therefore use assembly-derived phase, report errors inside and outside the benchmark regions separately, and stratify them by genomic context\cite{dwarshuis2024stratifications}.

Several limitations bound these conclusions. First, most data came from HG002, and the network was trained on HG002 nanopore runs with labels from the evaluation benchmark; genome-wide HG002 values include training chromosomes. On held-out chromosomes and in HG008, with weights, thresholds and graph constants fixed, \toolname kept its advantage over \whatshap and \hapcut, but GNN-based unphasing removed fewer switch errors than genome-wide in HG002 (Supplementary Note~7). The genome-wide HG002 values therefore overstate the effect to expect in a new sample. The held-out chromosomes informed the choice of label definition and earlier tuning of the graph constants, and the HG008 test covered nanopore SNVs at 10--30$\times$ against phase from an assembly built in part from the evaluated reads. Second, the HiFi results support transfer between platforms within HG002 only. Third, the other variant classes were evaluated less completely than SNVs. SV phase was scored for one SV caller against variants called from the assembly rather than a curated SV benchmark, 5mC evaluation measured phased fractions rather than accuracy, and neither class had supervised training labels. Methylation co-phasing was tested on autosomes of a male lymphoblastoid cell line\cite{fu2024methphaser}, and the method assumes diploidy and balanced allele fractions. Fourth, subsampling variability was not estimated above 20$\times$ for nanopore data or at any HiFi coverage. Evaluation in more individuals and tissues, for every variant class and against phase references independent of the evaluated reads, is needed to establish how far these results generalize.

Two directions follow from these results. First, \toolname currently makes a binary choice to phase or withhold. Its phasing entropy measures agreement among votes, not the strength of evidence, unlike the model-based confidence of \hapcut\cite{edge2017hapcut2} and statistical phasers\cite{hofmeister2023shapeit5}, and the GNN score is neither calibrated nor written to the output. Whether the GNN adds value over simpler rules, such as an entropy threshold or the confidence pruning of \hapcut at equal numbers of phased variants, remains to be tested; a calibrated per-variant confidence would let downstream methods weigh phase uncertainty directly. Second, pipelines for somatic and structural-variant analysis already use \longphase\cite{zheng2022clair3,zheng2026clairs,keskus2025severus,ont2025wfhuman}. Varying the phaser while fixing the somatic caller would show whether more accurate germline phase improves tumour-read assignment or whether withheld phase removes information that the caller needs. More broadly, these results show that genetic and epigenetic variation can be phased together from the same reads, with lower error rates than current phasers, when each assignment is weighed against all the evidence and uncertain assignments are withheld.
